% 67-50 ⑤(67쪽) — 보기 그래프: x=a 와 x=c 에서 x축을 지나고 x=b 에서 극소인 두 봉우리 개형.
% 그림 자체가 보기라 TikZ 로 다시 그림. 라벨은 원문 그대로 a · b · c · O · x · y 만.
\begin{tikzpicture}[x=0.46cm, y=0.46cm, line width=0.5pt, baseline=(current bounding box.center)]
  \path (-2.35,-1.75) rectangle (2.95,1.75);
  \draw[->, >=stealth, line width=0.7pt] (-2.00,0) -- (2.65,0) node[below=1pt, font=\small] {$x$};
  \draw[->, >=stealth, line width=0.7pt] (0,-1.55) -- (0,1.55) node[left=1pt, font=\small] {$y$};
  \draw[line width=0.6pt] plot[smooth, tension=0.7] coordinates {
    (-1.70,-1.30) (-1.55,-0.75) (-1.40,-0.35) (-1.20,0.00) (-1.00,0.35) (-0.80,0.60)
    (-0.60,0.72) (-0.40,0.68) (-0.20,0.50) (0.05,0.28) (0.30,0.12) (0.60,0.05)
    (0.85,0.25) (1.05,0.55) (1.20,0.78) (1.35,0.85) (1.50,0.75) (1.65,0.45)
    (1.75,0.20) (1.82,0.00) (1.95,-0.50) (2.10,-1.20)};
  \node[below=1pt, font=\small] at (-1.20,0) {$a$};
  \node[below left=-2pt and -3pt, font=\small] at (0,0) {$\pt{O}$};
  \node[below=1pt, font=\small] at (0.78,0) {$b$};
  \node[below=1pt, font=\small] at (1.90,0) {$c$};
\end{tikzpicture}
